% =====================================================================
\chapter{Aproximación universal y profundidad}
\label{tema:3}
% =====================================================================

¿Qué funciones puede representar un perceptrón multicapa? Este tema demuestra que una sola capa oculta basta, en principio, para aproximar cualquier función continua en un compacto, y explica por qué ese resultado no justifica por sí solo el uso de redes anchas y poco profundas. La demostración en una dimensión es constructiva y usa solo la continuidad uniforme; el caso general se enuncia según Cybenko~\cite{cybenko}, Hornik~\cite{hornik} y Leshno \emph{et al.}~\cite{leshno}. La discusión sobre la profundidad sigue la §6.4.1 de~\cite{gbc} y el artículo de Telgarsky~\cite{telgarsky}; la §5.1.1.4 de~\cite{d2l} y el apartado «Representational power» de~\cite{cs231n} ofrecen el mismo punto de vista.

\section{Redes ReLU y funciones lineales a trozos}

\begin{definicion}[Función lineal a trozos]
Una función continua $f\colon\R\to\R$ es \emph{lineal a trozos} con \emph{codos} $t_1<\dots<t_m$ si es afín en cada uno de los intervalos $(-\infty,t_1]$, $[t_1,t_2]$, \dots, $[t_m,\infty)$. Cada uno de esos $m+1$ intervalos es una \emph{pieza}.
\end{definicion}

\begin{proposicion}\label{prop:piezas}
Sea $f(t)=c+\sum_{j=1}^{H}a_j\ReLU(w_jt+b_j)$ una red de una capa oculta con $H$ unidades ReLU, entrada y salida escalares. Entonces $f$ es continua y lineal a trozos con a lo sumo $H$ codos, es decir, con a lo sumo $H+1$ piezas.
\end{proposicion}

\begin{lema}[Sombrero]\label{lem:sombrero}
La función $\Lambda(t)=\ReLU(t)-2\ReLU(t-1)+\ReLU(t-2)$ vale $0$ para $t\le 0$ y para $t\ge 2$, vale $t$ en $[0,1]$ y $2-t$ en $[1,2]$.
\end{lema}

Sumando sombreros desplazados y escalados se obtiene cualquier poligonal con nodos equiespaciados; el teorema siguiente usa una construcción equivalente y más económica.

\section{El teorema en una dimensión}

\begin{teorema}[Aproximación universal en $\R$ con ReLU]\label{teo:au1d}
Sea $f\colon[a,b]\to\R$ continua y sea $\varepsilon>0$. Existe una red $g(t)=c+\sum_{j=1}^{H}a_j\ReLU(t-\tau_j)$ con una capa oculta tal que $\abs{f(t)-g(t)}<\varepsilon$ para todo $t\in[a,b]$.
\end{teorema}

\addtocounter{definicion}{1}%

\section{El teorema general}

\begin{teorema}[Aproximación universal; Cybenko, Hornik, Leshno \emph{et al.}]\label{teo:au}
Sea $\phi\colon\R\to\R$ continua. Las funciones de la forma
\[
g(x)=\sum_{j=1}^{H}a_j\,\phi\bigl(x\cdot w_j+b_j\bigr),\qquad H\in\mathbb{N},\ a_j,b_j\in\R,\ w_j\in\R^{n},
\]
son densas en el espacio de funciones continuas sobre cualquier compacto $K\subset\R^{n}$ con la norma del supremo si y solo si $\phi$ no es un polinomio.
\end{teorema}

Cybenko~\cite{cybenko} lo demostró en 1989 para activaciones sigmoidales usando el teorema de Hahn--Banach y el de representación de Riesz; Hornik~\cite{hornik} lo extendió a activaciones acotadas y no constantes, y Leshno, Lin, Pinkus y Schocken~\cite{leshno} obtuvieron la caracterización anterior. La demostración del caso general excede el nivel de este curso. La necesidad de la condición sí es elemental:

\begin{proposicion}\label{prop:polinomio}
Si $\phi$ es un polinomio de grado $k$, toda red de la forma del \cref{teo:au} es un polinomio de grado a lo sumo $k$, y las redes no son densas en $C([0,1])$.
\end{proposicion}

\section{Qué no dice el teorema: el papel de la profundidad}

El \cref{teo:au} es un resultado de existencia. No acota el número $H$ de unidades necesarias, que para algunas funciones crece exponencialmente con la dimensión $n$ de la entrada~\cite[§6.4.1]{gbc}; no garantiza que el descenso de gradiente encuentre los pesos, y no dice nada sobre la generalización a datos no vistos. La profundidad aporta una ventaja que el teorema no recoge: componer capas permite reutilizar piezas, y algunas funciones que una red profunda representa con pocas unidades requieren exponencialmente muchas en una red poco profunda. El ejemplo siguiente lo muestra de forma exacta.

\begin{proposicion}[Diente de sierra]\label{prop:sierra}
Sea $g(t)=2\ReLU(t)-4\ReLU(t-\tfrac12)$, que en $[0,1]$ vale $2t$ en $[0,\tfrac12]$ y $2-2t$ en $[\tfrac12,1]$, y sea $g^{(k)}=g\circ\cdots\circ g$ ($k$ veces). Entonces:
\begin{enumerate}
  \item en $[0,1]$, $g^{(k)}$ es lineal a trozos con exactamente $2^{k}$ piezas de longitud $2^{-k}$, y en cada una recorre todo $[0,1]$;
  \item $g^{(k)}$ se calcula con una red de $k$ capas ocultas de 2 unidades ReLU cada una;
  \item toda red de una capa oculta que coincida con $g^{(k)}$ en $[0,1]$ necesita al menos $2^{k}-1$ unidades.
\end{enumerate}
\end{proposicion}

Con $k=20$, la red profunda usa 40 unidades y la poco profunda más de un millón. Telgarsky~\cite{telgarsky} demostró que la separación persiste aunque solo se pida \emph{aproximar}: hay funciones calculables por redes de profundidad del orden de $k^{3}$ y anchura constante que una red de profundidad $O(k)$ solo aproxima con del orden de $2^{k}$ unidades. Los LLM son profundos por esta razón: LLaMA-7B apila 32 bloques~\cite[§2]{touvron}, cada uno con un MLP de anchura moderada.

